\documentclass[11pt]{article}
\usepackage[T1]{fontenc}
\usepackage{lmodern,amsmath,amssymb,amsthm}
\usepackage[margin=1in]{geometry}
\setlength{\emergencystretch}{2em}
\usepackage[colorlinks=true,linkcolor=blue,citecolor=blue,urlcolor=blue]{hyperref}
\newtheorem{theorem}{Theorem}
\newtheorem{proposition}[theorem]{Proposition}
\newtheorem{lemma}[theorem]{Lemma}
\newcommand{\BM}{d_{\mathrm{BM}}}
\newcommand{\Cb}{C_b}
\newcommand{\Hb}{H_b}
\newcommand{\id}{\mathrm{id}}
\newcommand{\Z}{\mathcal Z}
\title{Ordinary Banach--Mazur rigidity of von Neumann algebras}
\author{Alec Kriebel\\\small ORCID: \href{https://orcid.org/0009-0001-9320-500X}{0009-0001-9320-500X}}
\date{October 7, 2026}
\hypersetup{pdftitle={Ordinary Banach-Mazur rigidity of von Neumann algebras},pdfauthor={Alec Kriebel},pdfsubject={Consequence of bounded Hochschild vanishing; arbitrary predual scope}}
\begin{document}
\maketitle
\begin{abstract}
Using OpenAI's all-algebra bounded Hochschild vanishing theorem, we derive
local ordinary complex Banach--Mazur rigidity for every von Neumann
algebra and its canonical predual. The threshold depends on the algebra;
the conclusion is Jordan $*$-isomorphism. Roydor announced the general
cohomological implication in 2020, while his published Theorem~1.2 assumes
separable predual. We give an explicit extension of his geometric argument
using deformation stability on two fixed source algebras and reconstruction
of odd finite type-I summands from full abelian projections. This avoids
separability, moving-corner perturbation constants, and enumeration of
infinite cardinal degrees. The cohomological input and geometric machinery
are attributed external results; no new solution of the underlying
cohomology problem or universal numerical threshold is claimed.
\end{abstract}

\section{Statement, inputs and attribution}
All Banach spaces and linear maps are complex. Define
\[
 \BM(X,Y)=\inf_T\|T\|\,\|T^{-1}\|,
\]
where $T$ ranges over bounded bijective complex-linear maps with bounded
inverse; the infimum is infinity when there is no such map. The literal
formula gives $\BM(0,0)=0$ and $\BM(0,X)=\infty$ for $X\ne0$.
For a von Neumann algebra $M$, write $M_*$ for its canonical predual, with
$M=(M_*)^*$ isometrically. A Jordan $*$-isomorphism is a bijective
complex-linear map preserving involution and $x\circ y=(xy+yx)/2$.

\begin{theorem}\label{thm:main}
For every complex von Neumann algebra $M$ there is $\varepsilon_M>0$
such that, for every complex von Neumann algebra $N$, either condition
\[
 \BM(M,N)<1+\varepsilon_M,
 \qquad \BM(M_*,N_*)<1+\varepsilon_M
\]
implies that $M$ and $N$ are Jordan $*$-isomorphic. With this threshold,
Jordan $*$-isomorphism, algebra-space linear isometry, predual linear
isometry, and the two displayed strict distance conditions are equivalent.
\end{theorem}

The essential new input is the following attributed theorem.
\begin{theorem}[OpenAI, family 295]\label{thm:input}
For every complex von Neumann algebra $V$ and every $k\ge2$,
every bounded complex $k$-linear Hochschild cocycle $V^k\to V$
has a bounded complex $(k-1)$-linear primitive. Thus $\Hb^k(V,V)=0$.
\end{theorem}
This is \cite[Theorem 1.1]{OAI295}, at the pinned repository version
\texttt{adc7f124} (full identifier in the pinned bibliography link). The mathematical proof
of that input has received independent source audits recorded in the
supplement. We rely on its manuscript proof, not a claim that our
follow-on result has been formalized. No fresh Lean kernel build of the
upstream declarations was reproduced.

Roydor's published reduction \cite[Theorem 1.2]{Roydor2022} assumes
\emph{separable predual}, $\Hb^2(M,M)=0$, and closed $B^3(M,M)$ in
$Z^3(M,M)$. Degree-three vanishing supplies the latter condition, since
$B^3=Z^3$ and $Z^3$ is a closed kernel. For separable-predual $M$,
Theorem~\ref{thm:main} is therefore an immediate consequence of that
published theorem and Theorem~\ref{thm:input}. Roydor's earlier author
slides already announce the unrestricted conditional formulation
\cite[Theorem 2, PDF pp.~28--29]{RoydorSlides}. Our contribution is the
explicit fixed-source extension below and the consequence of the newly
available cohomology input, not priority for that conditional formulation
or invention of its geometric machinery.

The supplied complete 26-page publisher-formatted Roydor version, dated
December 2020 and carrying the article's DOI, was read in full. Its
printed Theorem~1.2 and type-I argument must not be replaced by the
less restricted slide statement. We use only Corollary~2.4,
Proposition~2.13, Lemma~3.1 and Theorem~3.2 of that article, and its
approximate-order estimate, Lemma~2.5. Their proofs have no separability
hypothesis. In Theorem~3.2's proof, the off-diagonal relation in the
heading of Claim~4 should read $T(e_{ij})\simeq g_{ij}+h_{ji}$, as used
immediately above and in Claim~7. We do not use the article's
Theorem~1.1 or its compressed-map invocation in Lemma~4.3.

\section{An involutive form of multiplication stability}
For a complex $C^*$-algebra $A$, let $\Cb^k(A,A)$ be all bounded complex
$k$-linear maps, with the usual multilinear norm; $\Cb^0(A,A)=A$.
We use actual-image Hochschild cohomology, with
\begin{align*}
 (df)(a_1,\ldots,a_{k+1})={}&a_1f(a_2,\ldots,a_{k+1})\\
 &+\sum_{j=1}^k(-1)^j f(a_1,\ldots,a_ja_{j+1},\ldots,a_{k+1})\\
 &+(-1)^{k+1}f(a_1,\ldots,a_k)a_{k+1},
\end{align*}
and $(db)(a)=ab-ba$ in degree zero. Thus
$\Hb^k=\ker d/\operatorname{im}d$, with no closure in the denominator.
Neither normality nor complete boundedness is required. The differential
printed in \cite[Sec.~4.2.1]{Roydor2022} omits the last adjacent merge and
has an incorrect final sign; the standard differential above is the
intended one, agrees with family 295, and is used throughout this paper.

The following is an involutive version of the classical multiplication
stability mechanism of Johnson and Raeburn--Taylor
\cite{Johnson,RaeburnTaylor}; see also the symmetry-preserving iteration
in \cite[Proposition 3.2]{RicardRoydor}. We include the argument to make
its field, involution and fixed-algebra constants explicit.
\begin{lemma}\label{lem:deform}
Let $A$ be a fixed complex unital $C^*$-algebra with
$\Hb^2(A,A)=\Hb^3(A,A)=0$. There is $\delta_A>0$ such that every bounded
associative complex bilinear product $\mu$ on $A$ satisfying
\[
 \|\mu-m_A\|\le\delta_A,
 \qquad \mu(x,y)^*=\mu(y^*,x^*)
\]
admits an invertible bounded complex-linear $*$-preserving map $\Phi$
with $\|\Phi-\id\|<1/2$ and
$\Phi\mu(x,y)=\Phi(x)\Phi(y)$. If $\mu$ has the original unit, $\Phi$
preserves it.
\end{lemma}
\begin{proof}
Put $Z^k=\ker(d:\Cb^k\to\Cb^{k+1})$. Since $\Hb^3=0$, the map
$d:\Cb^2/Z^2\to Z^3$ is a bounded bijection of Banach spaces. Since
$\Hb^2=0$, $d:\Cb^1\to Z^2$ is onto. The open mapping theorem gives
constants $K,L\ge1$ such that each $f\in\Cb^2$ has $z\in Z^2$ with
$\|f-z\|\le K\|df\|$, and each $z\in Z^2$ has a primitive $h$ with
$\|h\|\le L\|z\|$. Choose $K,L$ strictly larger than the respective
quotient inverse bounds, to accommodate nonattained infima; for a zero
right side choose the exact zero-error solution. A bounded linear
selection is unnecessary.

Define conjugate-linear isometric involutions, for $n\ge1$, by
\[
 S_n f(a_1,\ldots,a_n)
 =(-1)^{(n-1)(n-2)/2}f(a_n^*,\ldots,a_1^*)^*.
\]
Reversing the differential's terms gives $dS_n=S_{n+1}d$. Averaging
$z$ with $S_2z$, and its primitive with $S_1h$, preserves the preceding
bounds when $f=S_2f$. The $S_1$-fixed maps preserve the involution;
$S_2$-fixed bilinear maps satisfy the product symmetry in the statement.

Write $\mu=m_A+\Delta$, $\delta=\|\Delta\|$. Associativity gives
\[
 d\Delta(x,y,z)=\Delta(\Delta(x,y),z)-\Delta(x,\Delta(y,z)),
 \qquad \|d\Delta\|\le2\delta^2.
\]
Choose symmetric $z,h$ with
\[
 dh=z,\quad \|\Delta-z\|\le2K\delta^2,\quad
 \|h\|\le L(\delta+2K\delta^2).
\]
For $g=\id+h$ and $b=g^{-1}$, transport the product to
$\nu(x,y)=g\mu(bx,by)$. Its exact defect is
\[
 (\nu-m_A)(x,y)
 =\bigl[(\Delta-dh)+h\Delta-m_A(h\cdot,h\cdot)\bigr](bx,by).
\]
If $\delta\le1/(2K)$ and $\delta\le1/(4L)$, then
$\|h\|\le2L\delta\le1/2$, $\|b\|\le2$, and
\[
 \|\nu-m_A\|\le D\delta^2,
 \qquad D=8K+8L+16L^2.
\]
Transport preserves associativity and the involution. Take
$\delta_A=\min\{1/(2K),1/(16L),1/(2D)\}$ and iterate, using the
original differential of the fixed algebra $A$ at every stage.
Then $\delta_n\le2^{-n}\delta_0$ and
$\sum_n\|h_n\|\le4L\delta_0\le1/4$.
The ordered products $\Phi_n=g_{n-1}\cdots g_0$ converge in operator
norm and satisfy
$\|\Phi-\id\|\le\exp(4L\delta_0)-1<1/2$.
Thus $\Phi$ is invertible and $*$-preserving. Passing to the limit in
$\Phi_n\mu(x,y)=\mu_n(\Phi_nx,\Phi_ny)$ proves the conjugacy.
An onto algebra isomorphism sends a unit to a unit, proving the last
assertion when both original products have unit $1$.
\end{proof}

\section{Rigidity on a fixed algebra whose unit halves}
Say that a von Neumann algebra \emph{halves} if its unit is the sum of two
orthogonal equivalent projections. Equivalently, its finite type-I
homogeneous summands have even degrees only; properly infinite summands
and type-II summands also halve. This is the projection-halving criterion
used in Roydor's proof of Theorem~3.2, with the standard structure theory
\cite{KR,Takesaki}. All infinite cardinal dimensions are included.

\begin{lemma}\label{lem:halve}
Let $V$ be a fixed von Neumann algebra which halves and satisfies
$\Hb^2(V,V)=\Hb^3(V,V)=0$. There is $\eta_V>0$ such that every von
Neumann algebra $W$ with $\BM(V,W)<1+\eta_V$ is Jordan $*$-isomorphic
to $V$.
\end{lemma}
\begin{proof}
For a sufficiently small distance excess, Roydor's Corollary~2.4 supplies
a unital self-adjoint linear isomorphism $T:V\to W$ with
$\max\{\|T\|,\|T^{-1}\|\}\le1+t$, where $t\to0$ with that excess.
His Theorem~3.2 and Lemma~3.1 provide central $p\in V,q\in W$ and a
block isomorphism
\[
 F=T_p^q\oplus T_{1-p}^{1-q}
\]
The first block nearly preserves multiplication; the second nearly
reverses it. The estimates are
\begin{align*}
 \|F^{-1}\|&\le1+988\sqrt t,\\
 \max\{\|(T_p^q)^\vee\|,\|(T_{1-p}^{1-q})^\wedge\|\}
 &\le3147585\sqrt t.
\end{align*}
The superscripts denote the ordinary bilinear multiplicative and
anti-multiplicative defects, respectively. Zero blocks are omitted.
Put $y\diamond z=qyz+(1-q)zy$ on $W$. Centrality of $q$ makes this
an associative unital $C^*$-product on the same normed $*$-space, namely
the original product on $qW$ and the opposite product on $(1-q)W$.
The pulled-back product
\[
 \mu(x,y)=F^{-1}(F(x)\diamond F(y))
\]
is associative, unital and involution-compatible, and
\[
 \|\mu-m_V\|\le(1+988\sqrt t)3147585\sqrt t\longrightarrow0.
\]
Apply Lemma~\ref{lem:deform} once to the fixed algebra $V$.
Then $F\Phi^{-1}$ is a $*$-isomorphism from $V$ to $(W,\diamond)$,
hence a Jordan $*$-isomorphism to $W$ with its original product.
The constants are attached to $V$, not the moving projection $p$.
\end{proof}

\section{Full carriers and odd finite type-I summands}
The central carrier $c(h)$ of a projection is the least central
projection dominating it. We need the following quantitative observation.
\begin{lemma}\label{lem:carrier}
Let $T:C\to D$ be a unital self-adjoint linear isomorphism of von
Neumann algebras with both norms at most $1+t$, for sufficiently small
$t$. If $c_C(h)=1$ and a projection $k\in D$ satisfies
$\|T(h)-k\|\le r(t)=140\sqrt t$, then $c_D(k)=1$.
\end{lemma}
\begin{proof}
Roydor's Proposition~2.13 gives a center $*$-isomorphism $\theta$
with $\|T(z)-\theta(z)\|\le r(t)$ for central projections $z$.
Put $w=c_D(k)$ and $z=\theta^{-1}(w)$. His Lemma~2.5 applied to
$T^{-1}$ and $k\le w$ yields
\[
 T^{-1}(k)\le T^{-1}(w)+\gamma(t)1,
 \qquad \gamma(t)=2(1+t)\sqrt{2t+t^2}.
\]
Both $\|h-T^{-1}(k)\|$ and $\|T^{-1}(w)-z\|$ are at most
$(1+t)r(t)$. Therefore
\[
 h-z\le\bigl[2(1+t)r(t)+\gamma(t)\bigr]1.
\]
Choose $t$ so that the bracket is less than one. Since $z$ is central,
$h(1-z)$ is a projection; compressing the inequality by $1-z$ forces
it to be zero. Thus $h\le z$, and $c_C(h)=1$ forces $z=1$, $w=1$.
The same argument applies to complements, since $T$ is unital.
\end{proof}

\begin{lemma}\label{lem:odd}
Let $O$ be a fixed von Neumann algebra consisting of finite homogeneous
type-$I_n$ central summands with odd $n\ge3$. Choose a projection $e$
which has rank $n-1$ on each such summand. Suppose
$E=eOe$ satisfies $\Hb^2(E,E)=\Hb^3(E,E)=0$. Then $O$ is locally
Jordan rigid for ordinary Banach--Mazur distance among von Neumann
algebras, without a separability assumption.
\end{lemma}
\begin{proof}
Both $e$ and $1-e$ have full central carrier; $(1-e)O(1-e)$ is
abelian, and the fixed algebra $E$ halves. These assertions follow by
writing each finite homogeneous summand as its arbitrary abelian center
tensored with $M_n$, choosing the diagonal rank-$n-1$ projection,
and taking the bounded central direct product.

Let $T:O\to D$ be a unital self-adjoint near isomorphism as in
Lemma~\ref{lem:halve}. Round $T(e)$ to a projection $q$ using Roydor's
Proposition~2.13. Then $1-q$ is equally close to $T(1-e)$.
Lemma~\ref{lem:carrier} gives $c_D(q)=c_D(1-q)=1$.
Roydor's Lemma~3.1 gives near isomorphisms from $E$ to $qDq$, and
from $(1-e)O(1-e)$ to $(1-q)D(1-q)$, with both distortions tending
to one. His Proposition~2.13 applied to the latter compression makes
$(1-q)D(1-q)$ abelian. Its full carrier implies that $D$ is type I.
Lemma~\ref{lem:halve}, applied to the single fixed algebra $E$, gives
an exact Jordan $*$-isomorphism $J:E\to qDq$.

For full projection corners, compression identifies their centers with
the center of the original algebra. The center restriction of $J$
therefore induces a $*$-isomorphism $\beta:\Z(O)\to\Z(D)$.
It preserves arbitrary existing joins: a Jordan $*$-isomorphism and
its inverse are positive order isomorphisms, hence normal.
Let $z_n$ be the source homogeneous degree-$n$ central projection and
put $w_n=\beta(z_n)$. The corner $w_nqDq$ is homogeneous of degree
$n-1$. This follows from the usual central homomorphism/anti-homomorphism
decomposition of a von Neumann Jordan $*$-isomorphism; either part
preserves finite homogeneous degree.

Consequently $w_nq$ is the sum of $n-1$ orthogonal equivalent abelian
projections, each with central carrier $w_n$ in $D$. The complementary
projection $w_n(1-q)$ is abelian and also has carrier $w_n$.
Abelian projections with the same central carrier are equivalent
\cite{KR,Takesaki}. Thus $w_n$ is the sum of $n$ orthogonal equivalent
abelian projections, and $w_nD$ is homogeneous of degree $n$.
Its center is identified by $\beta$ with that of $z_nO$, so finite
homogeneous matrix classification gives $z_nO\cong w_nD$ as von
Neumann algebras. Taking their bounded central direct product gives
$O\cong D$. The centers can have arbitrary density. No measurable
fiber decomposition or infinite-cardinal cancellation was used.

All constants were chosen for the fixed $E$, and the compression
bounds depend only on the initial norm excess. Hence one threshold
works for all $D$, all $q$ and all finite odd degrees simultaneously.
\end{proof}

\section{Completion of the algebra and predual conclusions}
\begin{proof}[Proof of Theorem~\ref{thm:main}]
For $M\ne0$, use its three fixed central summands
\[
 M=A\oplus P\oplus O,
\]
where $A$ is type $I_1$, $O$ contains the finite odd type-$I_n$
summands with $n\ge3$, and $P$ contains everything else: finite even
and all infinite type-I summands, and all type-II and type-III summands.
Then $P$ halves, and the projection $e$ in Lemma~\ref{lem:odd}
chooses the second fixed source algebra $E=eOe$.
Theorem~\ref{thm:input} applies to each of $P,E$, since both are von
Neumann algebras. This is why the proof uses universal vanishing;
we do not infer the noncentral corner's cohomology merely by applying
a central restriction to $M$.

For sufficiently small ordinary distance, normalize the initial Banach
isomorphism using Roydor's Corollary~2.4. Its Proposition~2.13 sends the
three central projections to an orthogonal partition of $1_N$, and
Lemma~3.1 gives near isomorphisms from the fixed three source summands
to the corresponding target summands, with uniform distortion bounds
converging to one. The commutative pieces are $*$-isomorphic by
Proposition~2.13; Lemmas~\ref{lem:halve} and \ref{lem:odd} handle
$P$ and $O$. Their finite direct sum is a Jordan $*$-isomorphism
$M\to N$. Omit zero summands. This includes arbitrary nonseparable
centers and arbitrary infinite type-I cardinal dimensions.

Choose the deformation thresholds once for the at most two nonzero
fixed algebras $P,E$. There are only finitely many subsequent center,
compression, carrier and halving inequalities; all their norm error
bounds tend to zero as the original distortion tends to one. Choosing
a small positive $\varepsilon_M$ to satisfy them proves the stated
$\forall M\,\exists\varepsilon_M\,\forall N$ quantifiers. No universal
constant is inferred. An infimum-attaining map is unnecessary: a
strict distance bound supplies some isomorphism with distortion
strictly below that bound.

For the predual implication, if $S:M_*\to N_*$ is a bounded
complex-linear isomorphism, its Banach adjoint $S^*:N\to M$ has inverse
$(S^{-1})^*$ and exactly the same two norms. Taking infima gives
\[
 \BM(M,N)\le\BM(M_*,N_*).
\]
Thus the same algebra threshold suffices. Exact Jordan
$*$-isomorphisms are normal order isomorphisms and induce isometric
canonical-predual maps. Conversely, predual isometries adjoint to
algebra-space isometries, and Kadison's isometry theorem gives Jordan
$*$-isomorphism; see also \cite[Theorem 1.1]{Sherman}.
This proves all equivalences. For $M=0$, either strict distance
inequality forces $N=0$, so any positive threshold works.
\end{proof}

\section{Prior scope and verification limits}
Ricard--Roydor's unconditional theorem concerns \emph{completely bounded}
Banach--Mazur distance \cite[Theorem A]{RicardRoydor}; it cannot be used
from small ordinary distance, since $\BM\le d_{cb}$ has the wrong
direction for that inference. Its predual observation and conditional
two-matrix-level version are also inherited \cite[Remarks 3.4 and
3.7]{RicardRoydor}. Roydor's special ordinary-stable classes include
those whose type-$\mathrm{II}_1$ summand has property $\Gamma$, a
Cartan MASA or McDuff absorption \cite[Corollary 1.3]{Roydor2022}.
Neither those cases nor exact predual duality are new contributions.

A fresh primary-source priority audit on October 7, 2026 checked the
literature, equivalent formulations, companion manuscripts and current
upstream version. No explicit duplicate full arbitrary-algebra consequence
and proof was identified in that bounded search. This does not establish
firstness. The general conditional formulation was already announced
by Roydor; the present note gives a fixed-source proof extension and
records a consequence of OpenAI's newly public vanishing theorem.
Opposite algebras explain the Jordan conclusion: ordinary Banach norm
geometry need not determine the orientation of the associative product.

Independent mathematical source audits found no substantive gap in the
pinned cohomology manuscript, and separate adversarial checks examined
the fixed-source deformation, carrier and reconstruction arguments.
Finite computations in the supplement are checks of proof mechanisms,
not proofs of the infinite-dimensional input. Static authentication of
the actual upstream Lean sources is distinguished from kernel verification:
no fresh compilation/import or axiom probe was reproduced, and this
Banach--Mazur consequence is not formalized here. Complete-package
reviews and publication verification are recorded separately in the
accompanying research records.

\paragraph{Disclosure.}
AI tools were used extensively in the research, drafting and verification.
Automated adversarial reviews are internal checks, not conventional human
peer review. This preprint has not undergone conventional human peer
review or refereeing.

\begin{thebibliography}{9}
\bibitem{Roydor2022}
J.~Roydor, \emph{Banach--Mazur stability of von Neumann algebras},
Journal of Topology and Analysis \textbf{14}(3) (2022), 767--792.
\href{https://doi.org/10.1142/S1793525321500151}{doi:10.1142/S1793525321500151}.
Published online December 11, 2020.
\bibitem{RoydorSlides}
J.~Roydor, \emph{Banach--Mazur stability of von Neumann algebras},
CIRM lecture slides, October 2020, Theorem 2, PDF pp.~28--29.
\href{https://www.cirm-math.fr/RepOrga/2169/Slides/Roydor_Slides.pdf}{Primary author slides}.
\bibitem{OAI295}
OpenAI, \emph{Vanishing of higher bounded Hochschild cohomology},
OpenAI Math Release preprint, manuscript dated September 23, 2026.
\href{https://github.com/openai/math/blob/adc7f1241b42e322a6451854ab7e4b4c146bf78a/preprints/Vanishing-of-higher-bounded-Hochschild-cohomology-September-23-2026/paper.pdf}{Pinned source and paper}.
Citation metadata supplied by the manuscript-specific README.
\bibitem{RicardRoydor}
\'{E}.~Ricard and J.~Roydor, \emph{A non-commutative Amir--Cambern
theorem for von Neumann algebras and nuclear $C^*$-algebras},
Journal of Functional Analysis \textbf{267}(4) (2014), 1121--1136,
\href{https://doi.org/10.1016/j.jfa.2014.05.018}{doi:10.1016/j.jfa.2014.05.018};
\href{https://arxiv.org/abs/1108.1970v2}{arXiv:1108.1970v2}.
\bibitem{Johnson}
B.~E.~Johnson, \emph{Perturbations of Banach algebras},
Proceedings of the London Mathematical Society (3) \textbf{34}
(1977), 439--458.
\bibitem{RaeburnTaylor}
I.~Raeburn and J.~L.~Taylor, \emph{Hochschild cohomology and
perturbations of Banach algebras}, Journal of Functional Analysis
\textbf{25} (1977), 258--266.
\href{https://doi.org/10.1016/0022-1236(77)90072-6}{doi:10.1016/0022-1236(77)90072-6}.
\bibitem{KR}
R.~V.~Kadison and J.~R.~Ringrose, \emph{Fundamentals of the Theory of
Operator Algebras}, Vol.~II, American Mathematical Society, 1997.
\bibitem{Takesaki}
M.~Takesaki, \emph{Theory of Operator Algebras I}, Springer, 2002.
\bibitem{Sherman}
D.~Sherman, \emph{Noncommutative $L^p$ structure encodes exactly Jordan
structure}, \href{https://arxiv.org/abs/math/0309365}{arXiv:math/0309365v2}
(2004), Theorem 1.1.
\end{thebibliography}
\end{document}
